\documentclass[12 pt, letterpaper]{hmcpset}
\usepackage{hw-preamble}
\usepackage[color = pink]{todonotes}

\name{Jayden Thadani}
\class{MATH 82 --- Differential Equation}
\assignment{Homework 1}
\duedate{}

\begin{document}

\begin{problem}[1:]
	Verify that $y_{1}(t) = e^{-t}$ and $y_{2}(t) = \sinh{t}$ both satisfy the differential equation $y'' - y = 0$. Is $y(t) = A y_{1}(t) + B y_{2}(t)$ where $A$ and $B$ are arbitrary constants, also a solution? Why or why not?
\end{problem}

\begin{solution}
	It's easy to see that
	\[
		\begin{split}
			y_{1}'' & = \frac{d^{2} y_{1}}{dt^{2}} \\
				& = \frac{d}{dt}(- e^{-t}) \\
				& = e^{-t}
		\end{split}
	\]
	and also that
	\[
		\begin{split}
			y_{2}'' & = \frac{d^{2} y_{2}}{d t^{2}} \\
				& = \frac{d}{dt} (\cosh{t}) \\
				& = \sinh{t}.
		\end{split}
	\]
	Thus, both of them satisfy the differential equation $y'' - y = 0$.

	Note also that the differential operator on the vector space of $\mathcal{C}^{2}$ functions by $D : y \mapsto y'' - y$ is a linear map (it's the difference of the twice-differentiation and identity operators). Thus, its null space --- exactly the solution space of the differential equation --- is a vector subspace containing all linear combinations of $y_{1}$ and $y_{2}$. That is, $A y_{1} + B y_{2}$ is a solution for $A, B \in \mathbb{R}$.
\end{solution}

\pagebreak

\begin{problem}[2:]
	Suppose that an object is moving in a straight line with constant acceleration $a \in \mathbb{R}$. Use properties of integration to show that the position of the object as a function of time $t$ is given by
	\[
		p(t) = \frac{at^{2}}{2} + v_{0} t + p_{0},
	\]
	where $v_{0}$ and $p_{0}$ denote the velocity and position at time $t = 0$.
\end{problem}

\begin{solution}
	We have that $p''(t) = a$, a constant for all $t \in \mathbb{R}$. Thus,
	\[
		\begin{split}
			p'(t) & = \int_{0}^{t} a \, dt + p'(0) \\
			      & = at + v_{0}.
		\end{split}
	\]
	Similarly,
	\[
		\begin{split}
			p(t) & = \int_{0}^{t} at + v_{0} \, dt + p(0) \\
			     & = \frac{at^{2}}{2} + v_{0} t + p_{0}
		\end{split}
	\]
	as desired.
\end{solution}

\pagebreak

\begin{problem}[3:]
	Verify that
	\[
		y(t) = e^{t^{2}} \int_{0}^{t} e^{-s^{2}} \, ds + e^{t^{2}}
	\]
	is a solution to the differential equation $y' - 2t y = 1$ with initial condition $y(0) = 1$.
\end{problem}

\begin{solution}
	It's clear that $y(0) = 0$. Now we compute $y'$.
	\[
		\begin{split}
			y'(t) & = 2t e^{t^{2}} + 2t e^{t^{2}} \int_{0}^{t} e^{- s^{2}} \, ds + e^{t^{2}} (e^{-t^{2}} - e^{0}) \\
			      & = 2t (e^{t^{2}} + e^{t^{2}} \int_{0}^{t} e^{- s^{2}} \, ds ) + 1 \\
			      & = 2t y + 1.
		\end{split}
	\] \todo[inline]{fix $e^{t^{2}}(e^{- t^{2}} - e^{0}) = 1 - e^{t^{2}} \neq 1$}
	
	Rearranging gives us $y' - 2t y = 1$ as desired
\end{solution}

\pagebreak

\begin{problem}[4:]
	Determine if $y(x) = x^{2} + 1$ is a solution to the initial value problem consisting of $4yy' = (y')^{3} - 3 y'' x$ and the initial condition $y(0) = 1$.
\end{problem}

\begin{solution}
	No, it's not. Just notice that the left hand side is always divisible by $4$ nut the right hand side is not. More rigorously, $4yy' = 8 x^{3} + 8 x$ and $(y')^{3} - 3 y'' x = 8 x^{3} - 6 x$ which are not equal as functions.
\end{solution}

\pagebreak

\begin{problem}[5:]
	Write $\cos{\theta}$ and $\sin{\theta}$ as linear combinations of $e^{-i \theta}$.
\end{problem}

\begin{solution}
	\[
		\boxed{
			\begin{aligned}
				\cos{\theta} & = \frac{e^{i \theta} + e^{- i \theta}}{2} \\
				\sin{\theta} & = \frac{e^{i \theta} - e^{- i \theta}}{2 i}
			\end{aligned}
		}	
	\]

	We know that $e^{i \theta} = \cos{\theta} + i \sin{\theta}$. Thus, $e^{- i \theta} = \cos{\theta} - i \sin{\theta}$. From there, it's easy to just notice that their sums and differences give us the sine and cosine functions like we want.
\end{solution}

\pagebreak

\begin{problem}[6:]
	Calculate the following, assuming $a, t, \omega$ are real numbers.
	\begin{enumerate}
		\item $\operatorname{Im}(e^{i \omega t})$.
		\item $\operatorname{Im}(e^{-i \omega t})$.
		\item $\operatorname{Re}(e^{(a + i \omega)t})$.
		\item $\operatorname{Im}(e^{(a + i \omega)t})$.
	\end{enumerate}
\end{problem}

\begin{solution}
	\begin{enumerate}
		\item \fbox{$\operatorname{Im} (e^{i \omega t}) = \sin{\omega t}$.}

			Just notice that $e^{i \omega t} = \cos{\omega t} + i \sin{\omega t}$
		\item \fbox{$\operatorname{Im}(e^{- i \omega t}) = - \sin{\omega t}$.}

			Similarly, $e^{ - i \omega t} = \cos{\omega t} - i \sin{\omega t}$.
		\item \fbox{$\operatorname{Re}(e^{a + i \omega t}) = e^{a} \cos{\omega t}$.}

			Here $e^{a + i \omega t} = e^{a}(\cos{\omega t} + i \sin{\omega t})$
		\item \fbox{$\operatorname{Im}(e^{a} \sin{\omega t})$.}

			Follows from the calculation above.
	\end{enumerate}
\end{solution}

\pagebreak

\begin{problem}[7:]
	For each of the following ordinary differential equations, indicate its order, whether it is linear or non-linear, whether it is autonomous or non-autonomous, and whether it is driven or non-driven.
	\begin{enumerate}
		\item $\dfrac{dg}{dx} + g^{3} = 0$.
		\item $\ddot{y}(t) + e^{t y(t)} = \cosh{t}$.
		\item $r^{2} R''(r) + r R'(r) - 5R(r) = 0$.
		\item $\ddot{\theta} + \sin{\theta} = 0$.
		\item $f''' = f' + xf + 4 \sin{x}$.
		\item $\frac{y'}{y} = 7$.
	\end{enumerate}
\end{problem}

\begin{solution}
	\begin{enumerate}
		\item is a linear, autonomous, non-driven differential equation.
		\item is a non-linear, non-autonomous, driven differential equation.
		\item is a linear,  non-autonomous, non-driven differential equation
		\item is a non-linear, autonomous, non-driven differential equation.
		\item is a linear, non-autonomous, driven differential equation.
		\item is a linear, autonomous, non-driven differential equation.
	\end{enumerate}
\end{solution}

\end{document}
